\scanpage{161}

On note alors que toutes les structures de graphe à homothéties $(S,A_i,Z)$, pour $i\in\deleted{\unsure{Z}}\pi$, satisfont la condition (GH1), et on a \unsure{les|en} termes de \inserted{la seule donnée} $\pi\subset\unsure{\mathcal P|\mathfrak P}(\mathcal P_2(S))$ une opération de $Z$ sur $\pi$, notée

\[
\left\{\begin{aligned}
 &\underset{\rotatebox{90}{$\in$}}{\lambda}\longmapsto\underset{\rotatebox{90}{$\in$}}{T_\lambda}
       \qquad \deleted{\unsure{\{s,t\}}}\qquad\text{définie par}\\[-2pt]
 &Z\qquad \unsure{\mathfrak S|\mathscr S}_\pi\\
 &\forall s,t\in S,\ i\in\pi,\qquad\{s,t\}\in A_i\Longleftrightarrow\{s,\lambda t\}\in A_{T_\lambda i}.
\end{aligned}\right.\tag*{\shortstack{\deleted{\unsure{(5)}}\\(6)}}
\]

\deleted{Soit \unsure{On} appelle \uline{multigraphe à opérateurs}}

Deux \deleted{$G_0=A\unsure{'|_0}$} graphes à homothéties \inserted{$(S,A_0,Z_0)(S,A_1,Z_1)$ sur le $=$ ens. de sommets} sont dits \uline{équivalents} \deleted{s'ils \unsure{=|se} peuvent se déduire l'un de l'autre par} \deleted{existe $\lambda\in Z_0$ tel que l'on ait} \deleted{$A_1=\deleted{\unsure{A}}T_\lambda(A_0)$, une opération} si on a $Z_0=Z_1$, soit $Z$, et s'il existe $\lambda\in Z$ tel qu'on ait $A_1=T_\lambda A_0$, \deleted{\unsure{donc}} ou encore si $A_1$ est $\in$ de l'ens $\pi$ construit précédemment, en termes de $(S,A_0,Z)$. Il est immédiat que c'est \inserted{bien} une \inserted{relation} \unsure{équivalence|d'équivalence}, plus restrictive que la relation d'isomorphie -- et $=$ que la relation d'isomorphie commutant à l'opération de $Z$.

La donnée, \deleted{pour} sur un ens. $S$, d'une classe d'équivalence $\pi$ en un sens de graphes à homothéties, sera appelée un \uline{multigraphe} \uline{à homothéties}. Posons alors

% archive sheet number: 160
